\documentclass[11pt]{article}% Shared preamble for the rigor documents (Lamport-style structured proofs).  \input this file after \documentclass.
\usepackage[T1]{fontenc}\usepackage{lmodern}\usepackage{microtype}
\usepackage[margin=1in]{geometry}
\usepackage{amsmath,amssymb,amsthm,mathtools}
\usepackage{xcolor}\usepackage{enumitem}\usepackage{booktabs}\usepackage{graphicx}\usepackage{hyperref}
\usepackage[most]{tcolorbox}
\definecolor{navy}{HTML}{17365D}\definecolor{teal}{HTML}{117A8B}\definecolor{warn}{HTML}{A33A2B}\definecolor{okgreen}{HTML}{2E7D4F}
\hypersetup{colorlinks=true,linkcolor=navy,citecolor=teal,urlcolor=teal}
\theoremstyle{plain}
\newtheorem{theorem}{Theorem}[section]\newtheorem{lemma}[theorem]{Lemma}\newtheorem{proposition}[theorem]{Proposition}\newtheorem{corollary}[theorem]{Corollary}
\theoremstyle{definition}\newtheorem{definition}[theorem]{Definition}\newtheorem{remark}[theorem]{Remark}\newtheorem{claim}[theorem]{Claim}\newtheorem{issue}[theorem]{Issue}
% ---------- Lamport structured proofs ----------
% \lstep{level}{number}{assertion}   -> indented "<level>number. assertion"
% \lproof                             -> "PROOF:" line (start of the sub-proof of the preceding step)
% \lby{justification}                 -> "BY ..." leaf justification for the preceding step
% \lqed{level}{number}                -> QED step of a sub-proof
% keywords: \lassume \lprove \lsuffices \lcase \ldefine \lpick \lnew
\newlength{\lindent}\setlength{\lindent}{1.7em}
\newcommand{\lstep}[3]{\par\smallskip\noindent\hspace*{\dimexpr\numexpr#1-1\relax\lindent\relax}{\color{navy}$\langle#1\rangle#2$.}\ #3}
\newcommand{\lproof}{\par\noindent\hspace*{\lindent}{\scshape Proof:}}
\newcommand{\lby}[1]{\par\noindent\hspace*{\lindent}{\scshape By}\ #1.}
\newcommand{\lqed}[2]{\par\smallskip\noindent\hspace*{\dimexpr\numexpr#1-1\relax\lindent\relax}{\color{navy}$\langle#1\rangle#2$.}\ {\scshape Q.E.D.}}
\newcommand{\lassume}{{\scshape Assume}\ }\newcommand{\lprove}{{\scshape Prove}\ }\newcommand{\lsuffices}{{\scshape Suffices}\ }
\newcommand{\lcase}{{\scshape Case}\ }\newcommand{\ldefine}{{\scshape Define}\ }\newcommand{\lpick}{{\scshape Pick}\ }\newcommand{\lnew}{{\scshape New}\ }
% ---------- verdict boxes ----------
\newtcolorbox{verdictok}{colback=okgreen!8,colframe=okgreen,title={\bfseries Verdict: verified},fonttitle=\small}
\newtcolorbox{verdictgap}{colback=orange!10,colframe=orange!80!black,title={\bfseries Verdict: gap / caveat},fonttitle=\small}
\newtcolorbox{verdictbreak}{colback=warn!10,colframe=warn,title={\bfseries Verdict: proof-breaking issue},fonttitle=\small}
\newtcolorbox{citedbox}[1]{colback=navy!5,colframe=navy!60,title={\bfseries Cited result: #1},fonttitle=\small,breakable}
\setlength{\parindent}{0pt}\setlength{\parskip}{4pt}\allowdisplaybreaks
\newcommand{\cA}{\mathcal A}\newcommand{\cF}{\mathcal F}\newcommand{\cM}{\mathfrak M}\newcommand{\cN}{\mathfrak N}

% extra calligraphic shortcuts (\providecommand: no clash if a document defines its own)
\providecommand{\cB}{\mathcal B}\providecommand{\cC}{\mathcal C}\providecommand{\cD}{\mathcal D}\providecommand{\cE}{\mathcal E}\providecommand{\cG}{\mathcal G}\providecommand{\cH}{\mathcal H}\providecommand{\cI}{\mathcal I}\providecommand{\cJ}{\mathcal J}\providecommand{\cK}{\mathcal K}\providecommand{\cL}{\mathcal L}\providecommand{\cO}{\mathcal O}\providecommand{\cP}{\mathcal P}\providecommand{\cQ}{\mathcal Q}\providecommand{\cR}{\mathcal R}\providecommand{\cS}{\mathcal S}\providecommand{\cT}{\mathcal T}\providecommand{\cU}{\mathcal U}\providecommand{\cV}{\mathcal V}\providecommand{\cW}{\mathcal W}\providecommand{\cX}{\mathcal X}\providecommand{\cY}{\mathcal Y}\providecommand{\cZ}{\mathcal Z}
\newcommand{\Fid}{\operatorname{F}}\newcommand{\Tr}{\operatorname{Tr}}\newcommand{\eps}{\varepsilon}\newcommand{\dd}{\,\mathrm d}

\usepackage{tabularx}
\newcommand{\Om}{\Omega}
\newcommand{\Hil}{\mathcal H}
\newcommand{\Hfin}{\Hil^{\mathrm{fin}}}
\newcommand{\Diff}{\mathrm{Diff}_+}
\newcommand{\ipp}[2]{\langle #1,#2\rangle}
\newcommand{\nrm}[1]{\lVert #1\rVert}
\newcommand{\abz}[1]{\lvert #1\rvert}
\newcommand{\gtot}{g_{\mathrm{tot}}}
\newcommand{\ktil}{\widetilde k}
\newcommand{\Ad}{\operatorname{Ad}}
\newcommand{\Exp}{\operatorname{Exp}}
\newcommand{\supp}{\operatorname{supp}}

\title{\color{navy}\bfseries Referee report QR2 on \texttt{rigor/implementation\_C11.tex}\\[0.3em]
\large Does the $C^{1,1}$ piecewise-M\"obius map really admit an implementer in \emph{every}
diffeomorphism covariant conformal net?}
\author{Referee QR2}\date{2026-09-08}

\begin{document}\maketitle
\vspace{-1.0em}

\begin{abstract}\noindent
Agent QB claims (Theorem~5.1 of the note) that hypotheses (H1)--(H3) of the Phase-2 brief hold for every
diffeomorphism covariant conformal net on $S^1$, with no restriction on $c$, no strong additivity and no
finiteness of $L_0$-multiplicities.  I re-read the two sources (Carpi--Weiner 2005 = CW05,
\texttt{refs/P1\_CarpiWeiner\_math-0407190.pdf}; Carpi--Del Vecchio--Iovieno--Tanimoto = CDIT,
\texttt{refs/CDIT\_1808.02384.pdf}) including the \emph{proofs} of the cited statements, checked every step of
the note's four proofs, re-derived the Fourier asymptotics and re-ran \texttt{rigor/qb\_regularity.py}.
\textbf{Verdict: Theorem~5.1 stands} --- (H1), (H2), (H3) are proved for every conformal net (CDIT \S4 axioms,
in particular on a separable Hilbert space), modulo one small, standard, repairable gap (the phase in
$U(\Exp f)=e^{iT(f)}$ must be a \emph{scalar} on a reducible $\Diff$-representation, and CDIT state that identity
only on the irreducible $\Hil(c,h)$).  The critical worry raised in my charge --- that CW05's essential
self-adjointness might go through Nelson's commutator theorem and hence need $\nrm{\gtot'}_{3/2}<\infty$ ---
is \emph{unfounded}: their proof uses the $\eps$-uniform bound $\nrm{[L_n,e^{-\eps L_0}]}\le\sqrt q\,\abz n^{3/2}$
and nothing else.  One genuine error, outside (H1)--(H3): the note's ``single missing analytic estimate'', the
\emph{form} bound with $\beta<1$, is not open --- it is \textbf{false}.  QB's own test vectors, weighted
optimally, force $\beta\ge1$; and the true endpoint bound ($\beta=1$, which I prove here in three lines from
CW05 Lemma~4.1) is useless for $\gtot$ because $\nrm{\gtot'}_1=\infty$ logarithmically.  The Nelson route to
$D(L_0)$-invariance is therefore closed, not open.
\end{abstract}

\tableofcontents

\section{Method, and what was actually checked}\label{sec:method}

I did not take the note's citedboxes on trust.  For each external input I opened the PDF at the page named in
the box, read the \emph{statement and its proof}, and checked that the hypotheses are the ones QB verifies:
CW05 pp.~11--14 (Lemma~4.1, Prop.~4.2, Prop.~4.3, Thm.~4.4, Prop.~4.5), CW05 pp.~15--18 (Lemma~4.6, Lemma~5.3,
Prop.~5.4), CDIT pp.~5--6 (\S2.1, Prop.~2.1, Prop.~2.2, Prop.~2.3), CDIT pp.~22--24 (net axioms (1)--(10),
Prop.~4.1, Prop.~4.2, Outlook).  Every transcription in the note's \S3 is faithful; two of them are sharper than
the printed source only in the direction QB claims (see \S\ref{sec:esa}).  All numerical claims were reproduced
by re-running \texttt{rigor/qb\_regularity.py} (4.7\,s, output in \S\ref{sec:numerics}).

\begin{center}\small
\begin{tabularx}{\textwidth}{>{\raggedright\arraybackslash}p{0.30\textwidth} >{\raggedright\arraybackslash}X c}\toprule
Ingredient of Theorem~5.1 & Finding of this report & Verdict\\\midrule
(1) $\gtot\in C^{1,1}$, one jump of $\gtot''$, $\hat G_n=\frac J{2\pi(in)^3}+O(n^{-4})$, $\nrm\gtot_{3/2}<\infty$
 & correct; asymptotics re-derived, constants and numerics reproduced (\S\ref{sec:reg}) & ok\\
(2) CW05 Thm.~4.4 $\Rightarrow$ $T(\gtot)$ e.s.a.\ on $\Hfin$; Stone
 & correct; the proof uses $\nrm{[L_n,e^{-\eps L_0}]}\le\sqrt q\abz n^{3/2}$, \emph{not} Nelson (\S\ref{sec:esa})
 & ok\\
(3) $\ktil_s=\Exp(s\gtot)$ is the compression map; $h_s\circ h_t=h_{s+t}$
 & correct; verified independently, incl.\ the generator $-x^2/L$ (\S\ref{sec:flow}) & ok\\
(4) covariance on $\cA(A)$ and $\cA(D)$: Trotter + localisation
 & correct; the decomposition is $\gtot=(\gtot-m)+m$, \emph{not} $g_A+g_D$ (\S\ref{sec:trotter}) & ok\\
(5) covariance on intervals straddling $0$: CDIT Prop.~4.1 method
 & correct; additivity/semicontinuity are automatic; one step needs a line (\S\ref{sec:straddle}) & ok (minor)\\
(6) $U(\Exp f)=e^{iT(f)}$ mod phase on a \emph{reducible} $\Diff$-representation
 & used essentially twice; CDIT state it on $\Hil(c,h)$ only (\S\ref{sec:phase}) & gap\\
(7) (H2): $\Om\in D(L_0)\subseteq D(\overline{T(\gtot)})$, Stone differentiability, signs
 & correct, signs agree with the brief and with Theorem~B (\S\ref{sec:H2}) & ok\\
(8) (H3) on all of $\cA(I)$, and $\omega_s(X^2)\to\omega(X^2)$
 & correct; ``$i\omega([T,X])$'' is formal but matches Theorem~B's $\varphi$ (\S\ref{sec:H3}) & ok (minor)\\
(9) ``form bound with $\beta<1$ is open, window $\tfrac12\le\beta<1$''
 & \textbf{false}: the same test family forces $\beta\ge1$; $\beta=1$ holds and is sharp (\S\ref{sec:form})
 & break\\
(10) numerics of \texttt{qb\_regularity.py}
 & reproduced to the digits quoted (\S\ref{sec:numerics}) & ok\\
\bottomrule
\end{tabularx}
\end{center}

\section{Ingredient 1: the regularity of $\gtot$ (note Thm.~2.2)}\label{sec:reg}

I re-derived (a)--(g) independently.  The one step where an error would have been fatal is (b), because the whole
construction rests on $\nrm\gtot_{3/2}<\infty$, i.e.\ on a decay \emph{strictly better} than $\abz n^{-5/2}$.

\lstep{1}{1}{$G$ is $C^1$ on $S^1$ and $C^\infty$ off $\theta=0$, with $G''(0^-)=0$, $G''(0^+)=-1/L$.}
\lby{$x=\tan(\theta/2)$, $G(\theta)=2\cos^2(\theta/2)g(x)$; on the $D$ side $g=-x^2/L$ gives
$G(\theta)=-\frac2L\sin^2(\theta/2)=-\frac{\theta^2}{2L}+O(\theta^4)$, on the $A$ side $G\equiv0$; hence
$G(0)=G'(0)=0$ and the second derivatives are $0$ and $-1/L$.  Note the exact circle field on $[0,L]$ is
$-\frac2L\sin^2(\theta/2)$, i.e.\ the M\"obius field $M$ of (C6), so the jump is $J=-1/L$ as claimed}
\lstep{1}{2}{$(in)^3\hat G_n\to J/2\pi$ and $\abz{\hat G_n}=\frac{\abz J}{2\pi\abz n^3}(1+O(\abz n^{-1}))$.}
\lby{$G'''=\{G'''\}+J\delta_0$ in $\mathcal D'(S^1)$ with $\{G'''\}$ piecewise smooth and bounded, hence BV;
$\widehat{\delta_0}=\frac1{2\pi}$ and $\widehat{\{G'''\}}_n=O(1/n)$ for BV; three integrations by parts on the
circle produce no boundary terms because $G,G',G''$ are integrable and $G,G'$ are continuous}
\lstep{1}{3}{$\nrm G_{3/2}<\infty$, $\nrm{G'}_{3/2}=\frac{2\abz J}\pi\sqrt N+O(\log N)$ truncated at $N$,
$G\in H^s\iff s<5/2$, $\nrm{T(G)\Om}^2=\frac c{12}\sum_{n\ge2}n(n^2-1)\abz{\hat G_n}^2<\infty$, and
$\nrm{L_0T(G)\Om}^2$ diverges like $\frac c{12}\frac{J^2}{4\pi^2}\log N$.}
\lby{$\langle1\rangle2$: $\sum\abz n^{-3}\abz n^{3/2}<\infty$; $\sum_{\abz n\le N}\abz n\abz{\hat G_n}\abz n^{3/2}
=2\frac{\abz J}{2\pi}\sum_{n\le N}n^{-1/2}+O(\log N)=\frac{2\abz J}\pi\sqrt N+O(\log N)$;
$\sum(1+n^2)^s\abz{\hat G_n}^2\asymp\sum n^{2s-6}$; $L_n\Om=0$ for $n\ge-1$ and
$\nrm{L_{-n}\Om}^2=\ipp\Om{[L_n,L_{-n}]\Om}=\frac c{12}n(n^2-1)$ give
$\sum n^{3}\abz{\hat G_n}^2\asymp\sum n^{-3}<\infty$ and, with the extra $n^2$,
$\sum n^{-1}$ with coefficient $(\abz J/2\pi)^2$}
\lqed{1}{4}\lby{$\langle1\rangle1$--$\langle1\rangle3$; all seven assertions of the note's Thm.~2.2 agree}

\begin{verdictok}
Theorem~2.2 of the note is correct as stated, including the constants.  The two facts that matter downstream are
$\nrm\gtot_{3/2}<\infty$ (which is \emph{all} that Carpi--Weiner need, \S\ref{sec:esa}) and
$\gtot\in\dot H^{3/2}$ (which is what makes $T(\gtot)\Om$ a vector).  The margin claimed in Remark~2.4 is real:
the class is $\abz{\hat G_n}\asymp n^{-3}$ against the two thresholds $n^{-5/2}$ (for $\nrm\cdot_{3/2}$) and
$n^{-2}$ (for $\dot H^{3/2}$).  I checked the $\Rightarrow$ direction of (e) as well: since $\gtot(0)=0$ the
singular point is a \emph{fixed point} of the flow, so $\ktil_s$ is smooth off $\theta=0$ for all $s$ and
$\ktil_s-\iota$ has exactly the same one-jump structure; $D^\sigma$ membership iff $\sigma<5/2$ follows, and
$\ktil_s\notin D^3$, so CDIT Cor.~3.6 is indeed inapplicable --- the note is not attacking a straw man.
\end{verdictok}

\section{Ingredient 2: essential self-adjointness --- the decisive point}\label{sec:esa}

My charge asked whether CW05 Thm.~4.4 gives essential self-adjointness under $\nrm f_{3/2}<\infty$ alone, or
whether its proof secretly uses Nelson's commutator theorem with $[L_0,T(f)]=iT(f')$ --- in which case
$\nrm{\gtot'}_{3/2}=\infty$ would destroy the application.  I read the proof (CW05 pp.~11--14).  It does not.

\begin{citedbox}{CW05 Lemma 4.1, Prop.~4.3, Thm.~4.4 (pp.~11--13): the actual mechanism}
\emph{Verbatim, Lemma 4.1 (p.~11).}  ``There exists a constant $r>0$ independent from $k,v_k,n$ (but dependent
on the value of the central charge $c$) such that $\nrm{L_nv_k}^2\le r^2(k^2+k\abz n^2+\abz n^3)\nrm{v_k}^2$,
where $v_k$ is an eigenvector of $L_0$ with eigenvalue $k$ and $n\in\mathbb Z$.''  Preceding sentence (p.~11):
``Although equation (2.8) \dots\ of Goodman and Wallach is stated for the irreducible case, the value of the
lowest weight is not involved at all \dots\ we have in general the following.''
\\
\emph{Verbatim, Prop.~4.3 (p.~12).}  ``There exists a constant $q>0$ independent of $\eps$ and $n$ such that
$\nrm{R_{n,\eps}}^2=\nrm{[L_n,e^{-\eps L_0}]}^2\le q\abz n^3$.''
\\
\emph{Verbatim, Thm.~4.4 (p.~13).}  ``If $a_n\in\mathbb C$ $(n\in\mathbb Z)$ is such that
$\sum_n\abz{a_n}\abz n^{3/2}<\infty$ then $A$ is closable and $\bar A=(A^+)^*$ \dots\ In particular, if
$a_n=\overline{a_{-n}}$ for all $n\in\mathbb Z$, then $A$ is essentially self-adjoint on $D^{\mathrm{fin}}$.''
\\
\emph{Proof mechanism (pp.~13--14), paraphrased faithfully.}  $\sum_n\nrm{a_nR_{n,\eps}}\le
q^{1/2}\sum_n\abz{a_n}\abz n^{3/2}<\infty$, so $R_{A,\eps}=[A,e^{-\eps L_0}]$ is \emph{bounded uniformly in
$\eps$}; $R_{A^+,\eps}\to0$ strongly on $D^{\mathrm{fin}}$ as $\eps\to0$, hence $R^*_{A,\eps}\to0$ strongly; then
for $x\in D(A^*)$, $A^+e^{-\eps L_0}x=e^{-\eps L_0}A^*x-R^*_{A,\eps}x\to A^*x$, so $A^*=\overline{A^+}$.
\\
\emph{Verdict.}  This is a mollifier/commutator argument with the \emph{bounded} operator $e^{-\eps L_0}$; the
only input is the one summability condition $\sum\abz{a_n}\abz n^{3/2}<\infty$, i.e.\ $\nrm f_{3/2}<\infty$
(for $\abz n\ge1$ the two conditions coincide, and $a_0$ contributes one bounded term).  \textbf{Nelson's
commutator theorem is not used}, no bound on $[L_0,T(f)]$ appears, and no smoothness of $f$ is used anywhere.
QB's claim (a) is therefore correct, and Prop.~5.1(iii)--(iv) of the note (``Nelson is unavailable and not
needed'') is right.  Two bonuses that the note under-uses: (i) CW05's $r$ is explicitly independent of the
lowest weight, and their \S4 is written for a \emph{general} positive-energy representation (see the quoted
sentence), so the note's Lemma~4.2 (reduction to irreducibles) is belt-and-braces rather than necessary ---
it is nevertheless correct, and it is needed if one insists on citing CDIT Prop.~2.2, which is stated on
$\Hil(c,h)$; (ii) Lemma~4.1 is the sharp input for \S\ref{sec:form} below.
\end{citedbox}

\begin{verdictok}
Step A1 of the note (essential self-adjointness of $T(\gtot)$ on $\Hfin$, hence a self-adjoint closure, hence
$U_s=e^{is\overline{T(\gtot)}}$ by Stone) is correct and correctly attributed.  The note's Lemma~4.2 (direct
sums, uniform $r(c)$) is proved correctly: the four assertions are summand-wise, the energy bound and the
resolvent convergence add in quadrature, and ``a symmetric extension of a self-adjoint operator equals it''
upgrades essential self-adjointness from the algebraic direct sum of the $\Hil_j^{\mathrm{fin}}$ to $\Hfin$.
\end{verdictok}

\section{Ingredient 3: is $\ktil_s$ really the flow of $\gtot$?}\label{sec:flow}

My charge asked to verify the one-parameter group property of the compression family and the identification of
the flow.  Both check out; I give the computations because the note only asserts them.

\lstep{1}{1}{$h_s\circ h_t=h_{s+t}$ for all $s,t$, with $h_s(x)=\frac{Lx}{L+sx}$.}
\lby{$h_s(h_t(x))=\dfrac{L\frac{Lx}{L+tx}}{L+s\frac{Lx}{L+tx}}
=\dfrac{L^2x}{L(L+tx)+sLx}=\dfrac{Lx}{L+(s+t)x}=h_{s+t}(x)$; in particular $h_0=\mathrm{id}$ and
$h_s^{-1}=h_{-s}$, so $\{h_s\}$ is a one-parameter group, as the note's $\langle2\rangle3$ needs}
\lstep{1}{2}{The generator is $m(x)=-x^2/L$, and $m$ is a globally smooth field on $S^1$.}
\lby{$\frac{d}{ds}h_s(x)\big|_{s=0}=-\frac{Lx\cdot x}{L^2}=-\frac{x^2}L$; more generally
$\frac{d}{ds}h_s(x)=-\frac{Lx^2}{(L+sx)^2}=m(h_s(x))$, so $s\mapsto h_s(x)$ is the integral curve of $m$.
In the circle chart $M(\theta)=\frac{2m}{1+x^2}=-\frac2L\sin^2(\theta/2)$, a trigonometric polynomial, so
$m\in\mathrm{Vect}(S^1)$ and $\Exp(sm)=h_s$ is M\"obius}
\lstep{1}{3}{$\ktil_s:=\Exp(s\gtot)$ exists, is a one-parameter group of homeomorphisms and equals $k_s$ on $I$.}
\lby{$\gtot\in C^1(S^1)$ hence Lipschitz, so Picard--Lindel\"of gives global existence and uniqueness on the
compact manifold $S^1$ and the flow property; uniqueness then forces $\ktil_s=\mathrm{id}$ on $(-a,0]$
(where $\gtot=0$) and $\ktil_s=h_s$ on $[0,L]$ (where $\gtot=m$ and, for $s\ge0$, $h_s([0,L])\subseteq[0,L]$
so the curve never leaves the region where $\gtot=m$)}
\lstep{1}{4}{$0$ is a fixed point of the flow; $\ktil_s$ is $C^{1,1}$, $C^\infty$ off $0$, with
$\ktil_s''(0^-)=0$, $\ktil_s''(0^+)=-2s/L$; and $\ktil_sI\subseteq I$ for $s\ge0$.}
\lby{$\gtot(0)=0$; $h_s'(x)=L^2/(L+sx)^2$, $h_s''(x)=-2sL^2/(L+sx)^3$ evaluated at $0$;
$\ktil_sI=(-a,\frac L{1+s})\subseteq I$}
\lqed{1}{5}\lby{$\langle1\rangle1$--$\langle1\rangle4$; the note's Prop.~2.1 is correct as stated}

\begin{verdictok}
The geometric layer is sound.  Two points that the note leaves implicit but that I verified because they are
load-bearing elsewhere: (i) the singular point is a fixed point, which is why $\ktil_s$ has \emph{one} bad point
for every $s$ (if $\gtot(0)\ne0$ the orbit of the kink would make $\ktil_s$ non-smooth at $\ktil_s(0)$ as well,
and Theorem~2.2(e) would still hold but the flow argument in \S\ref{sec:trotter} would need care); (ii) the
monotonicity $h_aD\subseteq D$ for $a\ge0$, on which the Trotter iteration depends, holds because $h_a$ fixes
$0$ and moves $L$ inward.  For $s<0$ this fails, which is why the note states its (ii) only for $s\ge0$; the
theorem's header ``for every $s\in\mathbb R$'' should be restricted there (cosmetic).
\end{verdictok}

\section{Ingredient 4: Trotter and localisation (note Thm.~4.4$\langle1\rangle$1--$\langle1\rangle$4)}
\label{sec:trotter}

My charge suggested the decomposition $\gtot=g_A+g_D$ into the restrictions to the two sides, and warned that
the pieces have kinks at $0$.  \textbf{That warning is correct and the note avoids the trap}: QB never uses
$g_A+g_D$.  It uses
\begin{equation}\label{eq:split}
 \gtot=(\gtot-m)+m,\qquad m=-x^2/L\ \text{globally smooth},\quad (\gtot-m)|_D=0 .
\end{equation}
This is essential, and worth recording explicitly:

\begin{verdictbreak}
\textbf{(For the brief, not for the note.)}  The decomposition proposed in the Phase-2 brief and in my charge,
$\gtot=g_A+g_D$ with $g_D:=\gtot\mathbf 1_{[0,L]}$, is \emph{unusable}, and for a worse reason than the
kink at $0$ that my charge anticipated: since $\gtot(L)=-L\ne0$, the truncation $g_D$ is \emph{discontinuous}
at $x=L$, so $\abz{\widehat{(g_D)}_n}\asymp\abz n^{-1}$ and $\nrm{g_D}_{3/2}=\infty$ (badly); the complement
$g_A=\gtot-g_D$ jumps there too.  Even the sharper split into $C^0$ pieces would fail: a field with a corner has
$\abz{\hat f_n}\asymp\abz n^{-2}$ and $\sum\abz n^{-2}\abz n^{3/2}=\infty$.  Neither $T(g_A)$ nor $T(g_D)$ is
then defined by CW05 at all, and Trotter would have nothing to work with.
The note's split \eqref{eq:split} keeps both summands in $\mathcal S^{3/2}$: $m$ is smooth, and $\gtot-m$
vanishes on $D$, equals $x^2/L$ on $(-a,0]$ and is $C^1$ with a single second-derivative jump at $0$
(jump $+1/L$), hence is again of the CDIT Prop.~2.3 type.
\end{verdictbreak}

With \eqref{eq:split} the note's four steps are:

\lstep{1}{1}{$\overline{T(\gtot)}=\overline{T(\gtot-m)+T(m)}$ with $\Hfin$ a common core, and Trotter applies.}
\lby{RS~I Thm.~VIII.31 needs $A=\overline{T(\gtot-m)}$, $B=\overline{T(m)}$ self-adjoint (CW05 Thm.~4.4, both
fields being real and of finite $3/2$-norm) and $A+B$ essentially self-adjoint on $D(A)\cap D(B)$.  The latter
holds: on $\Hfin\subseteq D(A)\cap D(B)$ the sum is $T(\gtot)$ term by term, which is e.s.a.\ there, and
$A+B$ on $D(A)\cap D(B)$ is a \emph{symmetric extension} of a self-adjoint operator, hence equals it.  This is
exactly the note's argument and it is correct}
\lstep{1}{2}{$\phi$ real, $\nrm\phi_{3/2}<\infty$, $\phi=0$ on an open interval $I_0$ $\Rightarrow$
$e^{it\overline{T(\phi)}}\in\cA(I_0)'$.}
\lby{CW05 Lemma~4.6 in the corrected reading (mollify: $\supp(\varrho_k*\phi)$ lies in the $\eps$-enlargement
$I_\eps$ of $\supp\phi\subseteq\overline{I_0'}$, and $\nrm{\varrho_k*\phi-\phi}_{3/2}\to0$ by dominated
convergence since $\abz{\hat\varrho_{k,n}}\le1$, $\hat\varrho_{k,n}\to1$); each $e^{itT(\phi_k)}$ is
$U(\Exp(t\phi_k))$ up to a phase, lies in $\cA(I_\eps')'=\cA(I_\eps)$ by axiom (10) and Haag duality
(CDIT (7), automatic); CW05 Prop.~4.5 gives $e^{itT(\phi_k)}\to e^{it\overline{T(\phi)}}$ strongly and von
Neumann algebras are SOT-closed; finally $\bigcap_{\eps>0}\cA(I_\eps)=\cA(I_0')=\cA(I_0)'$ by semicontinuity
(CDIT (9), automatic) plus Haag duality}
\lstep{1}{3}{$\Ad\bigl(e^{ia\overline{T(\gtot-m)}}e^{iaT(m)}\bigr)^n(x)=\Ad U(h_{na})(x)$ for $x\in\cA(D)$,
$a\ge0$; hence with $a=s/n$ and $\langle1\rangle1$, $\Ad U_s|_{\cA(D)}=\Ad U(h_s)|_{\cA(D)}$ for $s\ge0$.}
\lby{induction: $\Ad U(h_{ka})(x)\in\cA(h_{ka}D)\subseteq\cA(D)$ by $\langle1\rangle4$ of \S\ref{sec:flow}, and
on $\cA(D)$ the outer factor $e^{ia\overline{T(\gtot-m)}}\in\cA(D)'$ acts trivially by $\langle1\rangle2$ with
$I_0=D$; $U(h_{s/n})^n=U(h_s)$ modulo a phase since $\{h_s\}$ is a one-parameter group ($\langle1\rangle1$ of
\S\ref{sec:flow}) and $U$ restricted to $\mathrm{PSL}(2,\mathbb R)$ is a true representation; strong convergence
of unitaries passes to $\Ad$}
\lqed{1}{4}\lby{$\langle1\rangle1$--$\langle1\rangle3$; and $\Ad U_s|_{\cA(A)}=\mathrm{id}$ is
$\langle1\rangle2$ with $I_0=A_0\supseteq A$}

\begin{verdictok}
Correct.  Three details I checked separately because they are where such arguments usually fail.
\emph{(1) Trotter's hypotheses.}  RS~I Thm.~VIII.31 requires essential self-adjointness of $A+B$ on
$D(A)\cap D(B)$, not on the smaller $\Hfin$; the note's one-line upgrade (symmetric extension of a self-adjoint
operator) is valid and is needed.  \emph{(2) The order of the factors.}  With $V=e^{ia\overline{T(\gtot-m)}}$,
$W=e^{iaT(m)}$, $\Ad(VW)(x)=V(WxW^*)V^*$ and $WxW^*\in\cA(h_aD)\subseteq\cA(D)\subseteq\{V\}'$, so the induction
closes; had the M\"obius factor been on the outside the argument would need $\Ad V$ to preserve $\cA(D)$, which
is not given.  \emph{(3) $h_aD\subseteq D$ requires $a\ge0$}, which is why only the compressing direction is
obtained --- exactly the direction Phase~1 uses.  I also verified that CW05's own Prop.~5.4 is the same argument
(they have equality $\Ad e^{itT(g_p)}\cA(I_{(p,ip)})=\cA(I_{(p,ip)})$ because their M\"obius field vanishes at
both endpoints of the quarter-arc, whereas here one only has the inclusion; the note's induction handles that).
\end{verdictok}

\section{Ingredient 5: intervals straddling the touching point}\label{sec:straddle}

This is the step that CW05 Prop.~5.4 does \emph{not} give (they only need the four quarter-arcs on which the
map is M\"obius), and the note supplies it by running CDIT Prop.~4.1 with $e^{isT(g_n)}$ in place of
$U(\gamma_n)$.  My charge asks whether additivity is assumed and whether semicontinuity is used correctly.

\begin{citedbox}{CDIT \S4, p.~22: what is an axiom and what is automatic}
\emph{Verbatim.}  ``A M\"obius covariant net $(\cA,U,\Om)$ on $S^1$ is a triple \dots\ acting on a
\textbf{separable} complex Hilbert space $\Hil$ \dots\ (1) Isotony \dots\ (5) Vacuum vector.  \emph{With these
assumptions, the following automatically hold} [GF93, Theorem 2.19(ii)][FJ96, Section 3]: (6) Reeh--Schlieder
\dots\ (7) Haag duality: for every $I\in\mathcal I$, $\cA(I')=\cA(I)'$ \dots\ (8) Additivity: if
$\{I_\alpha\}$ is a covering of $I$ \dots\ then $\cA(I)\subset\bigvee_\alpha\cA(I_\alpha)$. (9)
Semicontinuity: if $I_n$ is a decreasing family of intervals and $I=(\bigcap_nI_n)^\circ$ then
$\cA(I)=\bigwedge_n\cA(I_n)$.''  Axiom (10) is diffeomorphism covariance with
$U(\gamma)x U(\gamma)^*=x$ for $x\in\cA(I)$, $\gamma\in\Diff(I')$.
\shot[\linewidth]{QB_CDIT_netaxioms.png}
\emph{Verdict on the usage.}  The note's claim that additivity, semicontinuity, Haag duality and
Reeh--Schlieder are \emph{consequences} of M\"obius covariance --- not extra hypotheses --- is exactly what the
source says, and it is the reason no strong additivity and no finite multiplicity enters.  One hypothesis that
does enter silently and should be stated in the note's (C4): \textbf{$\Hil$ is separable} (part of CDIT's
definition; it is also what makes the decomposition into irreducibles of CDIT Prop.~3.14 available for the
note's Lemma~4.2).  This is harmless --- every model is separable --- but it is a hypothesis.
\end{citedbox}

\lstep{1}{1}{$g_n\in C^\infty$ real with $\nrm{g_n-\gtot}_{3/2}\to0$; $\varphi^{(s)}_n:=\Exp(sg_n)$; then
$U(\varphi^{(s)}_n)=e^{isT(g_n)}\to U_s$ strongly and $\varphi^{(s)}_n\to\ktil_s$ uniformly.}
\lby{CW05 Lemma~4.6 with $I=S^1$; CW05 Prop.~4.5 (strong resolvent convergence, hence
$e^{itT(g_n)}\to e^{it\overline{T(\gtot)}}$ strongly, uniformly for $t$ in compacts, RS~I Thm.~VIII.21);
$\nrm{g_n-\gtot}_\infty\le\sum_k\abz{\hat g_{n,k}-\hat G_k}\le\nrm{g_n-\gtot}_{3/2}$ and Gronwall with the
Lipschitz constant of $\gtot$}
\lstep{1}{2}{For every proper interval $J$ and every $m$ large enough, $\bigcup_{k\ge m}\varphi^{(s)}_kJ$ is a
proper interval, and these decrease to $\overline{\ktil_sJ}$.}
\lby{uniform convergence of $\varphi^{(s)}_k$ and of their inverses moves the two endpoints of $J$ to within
$\delta_k\to0$ of those of $\ktil_sJ$; for $m$ with $\delta_m<\frac13\abz{\ktil_sJ}$ all the arcs
$\varphi^{(s)}_kJ$, $k\ge m$, contain the middle third of $\ktil_sJ$, hence pairwise intersect, so their union
is an arc, contained in the $\delta_m$-neighbourhood of $\ktil_sJ$ and therefore proper; the family decreases
in $m$; any point off $\overline{\ktil_sJ}$ is excluded for large $m$ and any point of $\ktil_sJ$ belongs to
$\varphi^{(s)}_kJ$ for all large $k$, so $\bigcap_m\overline{\bigcup_{k\ge m}\varphi^{(s)}_kJ}=\overline{\ktil_sJ}$,
whose interior is $\ktil_sJ$ because $\ktil_s$ is a homeomorphism}
\lstep{1}{3}{$U_s\cA(J)U_s^*\subseteq\cA(\ktil_sJ)$, and equality by applying the same to $U_{-s}$.}
\lby{for $x\in\cA(J)$ and $n\ge m$: $U(\varphi^{(s)}_n)xU(\varphi^{(s)}_n)^*\in\cA(\varphi^{(s)}_nJ)\subseteq
\cA(\bigcup_{k\ge m}\varphi^{(s)}_kJ)$ by axiom (10) and isotony (using $\langle1\rangle2$); the strong limit
$U_sxU_s^*$ stays in that algebra ($U_n\to U$ strongly with all unitary gives $U_n^*\to U^*$ strongly, hence
$U_nxU_n^*\to U_sxU_s^*$ strongly, and von Neumann algebras are SOT-closed); intersect over $m$ and use
semicontinuity (9) with $\langle1\rangle2$; $\ktil_{-s}=\ktil_s^{-1}$ by the group property}
\lqed{1}{4}\lby{$\langle1\rangle1$--$\langle1\rangle3$}

\begin{verdictok}
The straddling case is correctly handled, and this is the genuinely new part of the note relative to CW05.
Two remarks.  (i) The note's parenthesis ``for $m$ large the union is an interval, all its members being
contained in a small neighbourhood of $\ktil_sJ$ and each containing $\varphi^{(s)}_m J$-neighbours by uniform
convergence'' is the only place in the note where a topological statement is waved through; my
$\langle1\rangle2$ above is the missing half-page, and it goes through.  (ii) Additivity is used in the form
$\cA(\bigcup_{k\ge m}\varphi_kJ)=\bigvee_{k\ge m}\cA(\varphi_kJ)$ ($\subseteq$ by (8), $\supseteq$ by isotony)
and semicontinuity in the form (9); both are automatic, so no hypothesis beyond CDIT's axioms (1)--(10) is
smuggled in.  Note also that step $\langle1\rangle3$ does \emph{not} need $s\ge0$, so (H1) holds in both
directions even though the sharper statement $\Ad U_s|_{\cA(D)}=\Ad U(h_s)$ is proved only for $s\ge0$.
\end{verdictok}

\section{Ingredient 6: the phase in $U(\Exp f)=e^{iT(f)}$ on a reducible representation}\label{sec:phase}

This is the one place where I do not think the note's citation covers its use.

\begin{citedbox}{CDIT \S2.1, p.~5 (what is actually stated)}
\emph{Verbatim.}  ``\dots\ there is an irreducible, unitary, strongly continuous multiplier representation $U$
of $\widetilde{\Diff}(S^1)$ \dots\ such that $q(U(\Exp(f)))=q(e^{iT(f)})$ for all $f\in\mathrm{Vect}(S^1)$''.
The whole of CDIT \S2 is written on the irreducible lowest-weight space $\Hil(c,h)$; $q$ is the quotient of
$\mathcal U(\Hil(c,h))$ by the scalars.
\shot[0.98\linewidth]{QB_CDIT_prop21.png}
\emph{Use in the note.}  Twice, and both times essentially:
in Thm.~4.4$\langle2\rangle2$ (``$U(\Exp(t\phi_k))=e^{itT(\phi_k)}$ up to a phase, which is irrelevant since
$\cA(I_\eps)$ contains the phases'') to conclude $e^{itT(\phi_k)}\in\cA(I_\eps)$; and in
$\langle1\rangle5$--$\langle1\rangle6$ to identify $\Ad e^{isT(g_n)}$ with $\Ad U(\varphi^{(s)}_n)$.
\end{citedbox}

\begin{verdictgap}
\textbf{Gap (small, standard, repairable).}  In the vacuum representation of a conformal net, $U|_{\Diff}$ is in
general \emph{reducible}: by CDIT Prop.~3.14 (quoted in the note) it is $\bigoplus_jU^{(c,h_j)}$, and for a net
strictly larger than its Virasoro subnet the sum is infinite.  Applying the quoted identity summand by summand
gives only
\[
 U(\Exp(tf))=Z_t\,e^{it\overline{T(f)}},\qquad Z_t=\bigoplus\nolimits_je^{i\alpha_j t}\in\{L_n:n\in\mathbb Z\}',
\]
i.e.\ a unitary in the commutant of the Virasoro generators --- \emph{not} a scalar, unless the $\alpha_j$ can
be chosen independent of $j$.  If $Z_t$ is not scalar, then $e^{itT(\phi_k)}=Z_t^*U(\Exp(t\phi_k))$ need not lie
in $\cA(I_\eps)$ (the projections onto the $\Hil_j$ are not local), and the localisation step
Thm.~4.4$\langle1\rangle$2 --- which carries \emph{all three} of (i), (ii), (iii) --- would fail.
\emph{Repair.}  Fix a lift of $U|_{\Diff}$ to a unitary representation of the Virasoro group (the central
extension), which exists for positive-energy projective representations and is the setting of
[FH05, Prop.~5.1] cited by CDIT for the Bott--Virasoro multiplier.  In that lift the central element acts as the
\emph{scalar} $c$, which is common to all summands (CDIT Prop.~3.14, as the note itself stresses), and the
one-parameter subgroup $t\mapsto\Exp(tf)$ is represented by $e^{it\,\mathrm dU(f)}$ with
$\mathrm dU(f)=\overline{T(f)}$; the residual freedom is a single global $e^{i\alpha t}$.  Equivalently: CW05's
own \S5 (Prop.~5.4, Thm.~5.5) uses the identity in exactly this generality for an arbitrary 4-regular
conformal net, so the statement is available in the literature --- but the note cites CDIT \S2.1, which does
not state it.  \textbf{Action for QB: one sentence and one extra citation}, e.g.\ ``on a reducible
positive-energy representation the phase is scalar because $c$ is common to all irreducible summands and the
lift to the Virasoro group is global (FH05 Prop.~5.1; CW05 \S4 opening paragraph)''.  Alternatively, note that
$U(\Exp(tf))$ and $e^{itT(f)}$ have the same adjoint action on $\cA$ provided $Z_t\in U(\Diff)'\cap
\cA(S^1)$, which for an irreducible net is $\mathbb C\one$; that route needs $Z_t\in\cA(S^1)=B(\Hil)$ and
irreducibility of the net, and gives the conclusion for the \emph{adjoint} action, which is all the note uses.
\end{verdictgap}

I emphasise that this affects the \emph{write-up}, not the truth of Theorem~5.1: with the second route above
(irreducibility of the net, which follows from uniqueness of the vacuum, axiom (5)) one has
$Z_t\in\cA(S^1)'=\mathbb C\one$ as soon as $Z_t$ commutes with every local algebra, and $Z_t$ does commute with
$\cA(I)$ whenever both $U(\Exp(tf))$ and $e^{itT(f)}$ do, which is the case for $f$ supported in $I'$ once one
knows it for one of them; the argument is circular as stated and must be replaced by the Virasoro-group lift.

\section{Ingredient 7: (H2)}\label{sec:H2}

\lstep{1}{1}{$\Om\in D(L_0)\subseteq D(\overline{T(\gtot)})$ and
$\nrm{T(\gtot)\Om}^2=\frac c{12}\sum_{n\ge2}n(n^2-1)\abz{\hat G_n}^2<\infty$.}
\lby{$L_0\Om=0$ so $\Om\in D(L_0)$; the energy bound $\nrm{T(f)\xi}\le r\nrm f_{3/2}\nrm{(1+L_0)\xi}$
(CW05 Prop.~4.5, valid because $\nrm\gtot_{3/2}<\infty$); the norm formula from
$L_n\Om=0$ ($n\ge-1$) and $\nrm{L_{-n}\Om}^2=\frac c{12}n(n^2-1)$}
\lstep{1}{2}{$\nrm{e^{isT}\xi-\xi-isT\xi}=o(s)$ for self-adjoint $T$ and $\xi\in D(T)$; hence
$U_s\Om=\Om+is\,T(\gtot)\Om+o(s)$.}
\lby{spectral theorem and dominated convergence with dominating function $4\lambda^2\in L^1(\mu_\xi)$; the
sign is fixed by $U_s=e^{is\overline{T(\gtot)}}$, matching (H2) of the brief and Definition~4.3 of
\texttt{universality\_theorem\_B.tex} ($\Psi_s=U_s\Om$, $\nrm{\Psi_s-\Om+is\,T(\gtot)\Om}=o(s)$ there is the
same statement with the opposite sign convention for $\Psi_s$; see the check below)}
\lqed{1}{3}\lby{$\langle1\rangle1$, $\langle1\rangle2$}

\begin{verdictok}
(H2) is correct, and the derivative is correctly identified.  Sign check, since the charge asks for it: with
$U_s=e^{is\overline{T(\gtot)}}$, $\frac d{ds}U_s\Om|_{s=0}=i\,\overline{T(\gtot)}\Om$, which is the brief's
(H2).  Theorem~B uses $a:=T(\gtot)\Om$ and $b=i\dot\Psi_0=a$ with $\Psi_s=U_s^*\Om$ (its Prop.~6.2
$\langle1\rangle2$), i.e.\ $\dot\Psi_0=-iT(\gtot)\Om$ --- consistent, because Theorem~B's curve is the
\emph{recovered} vector $U_s^*\Om$.  Both conventions appear in the two documents and they agree.
The domain point is also right: differentiability of $s\mapsto U_s\Om$ needs $\Om\in D(\overline{T(\gtot)})$,
the domain of the \emph{closure}, which is what the energy bound provides ($D(L_0)\subseteq
D(\overline{T(\gtot)})$, CW05 Prop.~4.5); $\Om\in\Hfin$ would give the same, but the energy bound is the
cleaner route and is the one used.
\end{verdictok}

\section{Ingredient 8: (H3), and what ``$i\omega([T,X])$'' means}\label{sec:H3}

The note proves (H3) on \emph{all} of $\cA(I)$, not on a dense subalgebra.  I verified the expansion.

\lstep{1}{1}{$\omega_s(X)=\ipp{\xi_s}{X\xi_s}$ with $\xi_s=U_s^*\Om=\Om-is\,T(\gtot)\Om+r_s$, $\nrm{r_s}=o(s)$.}
\lby{$\omega_s=\omega\circ\Ad U_s$, so $\omega_s(X)=\ipp\Om{U_sXU_s^*\Om}=\ipp{U_s^*\Om}{XU_s^*\Om}$;
\S\ref{sec:H2}$\langle1\rangle2$ with $s\mapsto-s$}
\lstep{1}{2}{$\omega_s(X)=\omega(X)+s\varphi(X)+o(s)$ uniformly on bounded sets, with
$\varphi(X)=i\bigl(\ipp{T(\gtot)\Om}{X\Om}-\ipp\Om{XT(\gtot)\Om}\bigr)$.}
\lby{expand the sesquilinear form; the first-order terms are $\overline{(-is)}\ipp{T\Om}{X\Om}+(-is)\ipp\Om{XT\Om}
=is\ipp{T\Om}{X\Om}-is\ipp\Om{XT\Om}$; the remainder is bounded by
$\nrm X(2\nrm{r_s}(1+s\nrm{T\Om})+\nrm{r_s}^2+s^2\nrm{T\Om}^2)$}
\lstep{1}{3}{$\omega_s(X^2)\to\omega(X^2)$, and $\varphi(K)=-2\operatorname{Im}\ipp{T(\gtot)\Om}{K\Om}$ for
$K=K^*$.}
\lby{$s\mapsto\xi_s$ is norm continuous and $X^2$ is bounded; for $K=K^*$,
$\ipp\Om{KT\Om}=\overline{\ipp{T\Om}{K\Om}}$, so $\varphi(K)=i(z-\bar z)=-2\operatorname{Im}z$ with
$z=\ipp{T(\gtot)\Om}{K\Om}$}
\lqed{1}{4}\lby{$\langle1\rangle1$--$\langle1\rangle3$}

\begin{verdictok}
(H3) holds on all of $\cA(I)$, with the $o(s)$ uniform on norm-bounded sets --- strictly more than the brief
asks.  Two caveats, both cosmetic.
\emph{(i) The commutator is formal.}  The note writes
$\varphi(X)=i\bigl(\ipp{T(\gtot)\Om}{X\Om}-\ipp\Om{XT(\gtot)\Om}\bigr)=i\,\omega([T(\gtot),X])$.  The second
expression presupposes $X\Om\in D(\overline{T(\gtot)})$, which is \emph{not} known for a general $X\in\cA(I)$
(indeed $T(\gtot)$ is not affiliated to $\cA(I)$, and $\gtot$ is not supported in $I$).  The rigorous object is
the first expression.  This is harmless because it is exactly the object Theorem~B consumes:
\texttt{universality\_theorem\_B.tex} Definition~6.1 sets
$\varphi(K)=i\,\omega([T(\gtot),K])=-2\operatorname{Im}\ipp a{K\Om}$ with $a=T(\gtot)\Om$, i.e.\ it uses the
sesquilinear form, and $g_B$ depends on $a$ only through $\varphi$ (its Lemma~6.3).  Recommend: define
$\varphi$ by the sesquilinear expression and call $i\omega([T,X])$ a notation.
\emph{(ii)} Theorem~B's Prop.~6.4 needs $\varphi$ and the second-order control on a $*$-subalgebra closed under
squares; getting (H3) on all of $\cA(I)$ makes that step vacuous, which removes the MINOR finding recorded by
QA against \texttt{lemma\_second\_variation.tex} in the case at hand.
\end{verdictok}

\section{The claim I disagree with: the ``missing form estimate'' is false, not open}\label{sec:form}

The note's Remark~5.2 and the last row of its summary table assert:
\begin{quote}\itshape
``$\abz{\ipp\psi{T(f)\psi}}\le C_\beta\nrm f_\beta\nrm{(1+L_0)^{1/2}\psi}^2$ with $\beta<1$ \dots\ the operator
version is false for $\beta<3/2$ \dots\ for the form version the same test vectors give a weaker restriction:
with $\psi=\Om+\nrm{L_{-n}\Om}^{-1}L_{-n}\Om$ \dots\ \eqref{eq:testv} forces $\beta\ge1/2$ only.  \textbf{The
window $\tfrac12\le\beta<1$ is open}, and closing it (for general $c$) is the single missing analytic estimate
in this circle of ideas.''
\end{quote}
The test family is the right one; it was simply not optimised.  Put
\begin{equation}\label{eq:testv}
 f=2\cos(n\theta)\ (n\ge2),\qquad \psi_t=\Om+t\,e_n,\quad e_n:=\nrm{L_{-n}\Om}^{-1}L_{-n}\Om,\ t>0 .
\end{equation}

\lstep{1}{1}{$\ipp{\psi_t}{T(f)\psi_t}=2t\nrm{L_{-n}\Om}$ and
$\nrm{(1+L_0)^{1/2}\psi_t}^2=1+(1+n)t^2$.}
\lby{$T(f)=L_n+L_{-n}$; $T(f)\Om=L_{-n}\Om=\nrm{L_{-n}\Om}e_n$ and
$L_ne_n=\nrm{L_{-n}\Om}^{-1}[L_n,L_{-n}]\Om=\nrm{L_{-n}\Om}\Om$ (using $L_0\Om=0$, $L_n\Om=0$,
$[L_n,L_{-n}]=2nL_0+\frac c{12}n(n^2-1)$ and $\nrm{L_{-n}\Om}^2=\frac c{12}n(n^2-1)$), while $L_{-n}e_n$ sits at
energy $2n$ and is orthogonal to $\Om$ and $e_n$; the four terms are $0+t\nrm{L_{-n}\Om}+t\nrm{L_{-n}\Om}+0$.
For the energy: $L_0\Om=0$, $L_0e_n=ne_n$}
\lstep{1}{2}{$\displaystyle\sup_{t>0}\frac{\ipp{\psi_t}{T(f)\psi_t}}{\nrm{(1+L_0)^{1/2}\psi_t}^2}
=\frac{\nrm{L_{-n}\Om}}{\sqrt{1+n}}=\Bigl(\frac c{12}\Bigr)^{1/2}\frac{\sqrt{n(n^2-1)}}{\sqrt{1+n}}
\ \sim\ \Bigl(\frac c{12}\Bigr)^{1/2}n .}
\lby{$\max_{t>0}\frac{2tB}{1+At^2}=\frac B{\sqrt A}$ at $t=A^{-1/2}$, with $A=1+n$, $B=\nrm{L_{-n}\Om}$}
\lstep{1}{3}{Since $\nrm f_\beta=2(1+n^\beta)$, the form bound with exponent $\beta$ forces $\beta\ge1$.}
\lby{$\langle1\rangle2$ gives $(c/12)^{1/2}n(1+o(1))\le 2C_\beta(1+n^\beta)$ for all $n$}
\lstep{1}{4}{Conversely the endpoint $\beta=1$ \emph{holds}:
$\abz{\ipp\psi{T(f)\psi}}\le C\nrm f_1\nrm{(1+L_0)^{1/2}\psi}^2$ for all $\psi\in\Hfin$, $C=C(c)$.}
\lby{CW05 Lemma~4.1: $\nrm{L_nv_k}\le r(k^2+k\abz n^2+\abz n^3)^{1/2}\nrm{v_k}$ for an $L_0$-eigenvector $v_k$.
Splitting $\psi=\sum_k\psi_k$ and using $\abz{\ipp\psi{L_n\psi}}\le\sum_k\nrm{\psi_{k+n}}\,\nrm{L_n\psi_k}$
together with the elementary inequality
$(k^2+k\abz n^2+\abz n^3)^{1/2}\le3\abz n\,(1+k)^{1/2}(1+k+\abz n)^{1/2}$ (check the three regimes $k=0$,
$k\asymp\abz n$, $k\gg\abz n$) and Cauchy--Schwarz gives
$\abz{\ipp\psi{L_n\psi}}\le3r\abz n\nrm{(1+L_0)^{1/2}\psi}^2$; sum against $\abz{\hat f_n}$}
\lstep{1}{5}{$\nrm{\gtot'}_1=\sum_n\abz n\abz{\hat G_n}(1+\abz n)=\infty$ (logarithmically).}
\lby{Theorem~2.2(b) of the note: $\abz{\hat G_n}\sim\frac{\abz J}{2\pi\abz n^3}$, so the summand is
$\sim\frac{\abz J}{2\pi}\abz n^{-1}$}
\lqed{1}{6}\lby{$\langle1\rangle3$--$\langle1\rangle5$}

\begin{verdictbreak}
\textbf{Proof-breaking (for Remark~5.2 and the last row of the note's summary table; \emph{not} for
Theorem~5.1).}  The window ``$\tfrac12\le\beta<1$ is open'' is empty.  On the scale $\nrm\cdot_\beta$ the form
bound holds \emph{exactly} at $\beta\ge1$: it is true for $\beta=1$ ($\langle1\rangle4$, three lines from CW05
Lemma~4.1) and false for every $\beta<1$ ($\langle1\rangle3$).  QB's weaker restriction $\beta\ge1/2$ comes from
evaluating the quotient at the single unnormalised vector $\Om+e_n$; the supremum over the one-parameter family
$\Om+te_n$ is larger by $\sqrt n$, and that supremum is what a bound with a constant must dominate.
\emph{Consequence for the programme.}  Since $\nrm{\gtot'}_\beta<\infty$ requires $\beta<1$ and the estimate
requires $\beta\ge1$ --- and at the borderline $\nrm{\gtot'}_1=\infty$ logarithmically ($\langle1\rangle5$) ---
the route ``form bound $\Rightarrow$ Nelson/Glimm--Jaffe $\Rightarrow$ $U_s$ preserves $D(L_0)$ /
$\ipp{U_s^*\Om}{L_0U_s^*\Om}<\infty$ for fixed $s$'' is \textbf{closed}, not open, on this scale.  It misses by
exactly one logarithm, which is the same logarithm as in Theorem~2.2(g) ($T(\gtot)\Om\in D(L_0^{1/2})\setminus
D(L_0)$) --- consistent, and in fact explained by it: if $U_s$ preserved $D(L_0)$ one would expect
$\overline{T(\gtot)}\Om\in D(L_0)$, which is false.  Anything of this kind must use a genuinely different
(e.g.\ $L^2$-based, $\dot H^{1}$- or $\dot H^{1/2}$-type, or $f$-positivity/QEI) norm on the right-hand side, or
give up on $\ipp{U_s^*\Om}{L_0U_s^*\Om}<\infty$.  Recommended rewrite of the note's last paragraph: ``the form
bound holds with $\beta=1$ and this is sharp; $\nrm{\gtot'}_1$ diverges logarithmically, so the commutator route
is unavailable and Route~B (a Schwarzian energy bound, \S6) is the only remaining source of the fixed-$s$ energy
statement''.
\end{verdictbreak}

The note's other two statements in that neighbourhood are correct: the \emph{operator} bound
$\nrm{T(f)\xi}\le C\nrm f_\beta\nrm{(1+L_0)\xi}$ is indeed false for $\beta<3/2$ (take $\psi=\Om$,
$\nrm{T(f)\Om}=\nrm{L_{-n}\Om}\asymp n^{3/2}$, $\nrm{(1+L_0)\Om}=1$), and Nelson's commutator theorem is indeed
unavailable.  Note in passing a factor-4 slip in the note's Prop.~5.1(ii): the displayed
$\nrm{T(f)\Om}=(\frac c3n(n^2-1))^{1/2}$ should be $(\frac c{12}n(n^2-1))^{1/2}$, which is what its own proof
step $\langle1\rangle2$ writes; the conclusion $\beta\ge3/2$ is unaffected.

\section{Minor errors and presentational points}\label{sec:minor}

\begin{enumerate}[leftmargin=1.6em,label=\textup{(M\arabic*)}]
\item \emph{Prop.~5.1(ii)}: $\nrm{T(2\cos n\theta)\Om}^2=\frac c{12}n(n^2-1)$, not $\frac c3n(n^2-1)$
(the proof has it right).  Harmless.
\item \emph{Prop.~5.1(iii) and Remark~5.2}: $[N,T(\gtot)]=-iT(\gtot')$ should be $+iT(\gtot')$, since
$[L_0,L_n]=-nL_n$ and $\widehat{(g')}_n=in\hat g_n$ give $[L_0,T(g)]=-\sum n\hat g_nL_n=iT(g')$.  The same sign
slip is in \texttt{phase2\_brief.md}.  Only the modulus is used.
\item \emph{Theorem~4.4 header}: ``For every $s\in\mathbb R$'' followed by item (ii) ``for $0\le s$''.  Item (ii)
is genuinely restricted to $s\ge0$ (it uses $h_sD\subseteq D$); (i) and (iii) hold for all $s$.
\item \emph{Theorem~4.4$\langle1\rangle$7} (independence of the extension) reuses the Trotter iteration with
$\widehat U_a$ in place of $U(h_a)$; this needs $\Ad\widehat U_a(\cA(I))\subseteq\cA(I)$, which the note supplies
from (iii) and $\ktil_aI\subseteq I$, hence again only for $a\ge0$.  Correct, but the restriction should be
stated in the claim (as it is in the brief: ``$\omega_s$ \dots\ since $\ktil_s=k_s$ on $I$'').
\item \emph{Verdict box after Theorem~2.2}: the numerical quantities quoted there are two-sided Fourier sums
($\sum_{n\ne0}$) whereas the theorem's formulas are one-sided ($\sum_{n\ge2}$).  The conversions used
($4.076764/24=0.169865$, $38.945685/24=1.622737$, $\frac{J^2}{2\pi^2}\log16=0.14049$ for the two-sided
$L_0^2$-moment) are all consistent, but the reader has to reconstruct them; say so once.
\item \emph{(C4)}: the separability of $\Hil$ is part of CDIT's definition of a M\"obius covariant net and is
used (via Prop.~3.14) in Lemma~4.2.  State it.
\item \emph{\S5, class (a)}: ``$\cA(I)$ has trivial relative commutant in itself up to phases'' is not a
statement I can parse; if the point is that two implementers of the same automorphism of $\cA(I)$ differ by a
unitary in $\cA(I)'\cap\ldots$, it should be spelled out or deleted.  Nothing depends on it.
\item \emph{\S3, citedbox on CW05 Lemma~4.6}: the ``compact support'' caveat is correctly identified and
correctly repaired by the $\eps$-enlargement plus semicontinuity.  Good catch, and it is a real gap in CW05's
own Prop.~5.4 as printed.
\end{enumerate}

\section{Numerics: \texttt{rigor/qb\_regularity.py} re-run}\label{sec:numerics}

Re-run verbatim (4.7\,s, $N=2^{22}$, $L=1$, $M=3$, representative (i)).  Reproduced values, against the note:
\begin{center}\small
\begin{tabularx}{\textwidth}{>{\raggedright\arraybackslash}p{0.33\textwidth} >{\raggedright\arraybackslash}X c}\toprule
Claim in the note & Re-run & \\\midrule
(b) $n^3\abz{\hat G_n}=0.159155$ at $n=2048,8192$ vs $\abz J/2\pi=0.1591549$
 & $0.159155$, $0.159155$ (and $0.159113$ at $n=32768$: grid error) & ok\\
(c) $\nrm G_{3/2}=18.7040$ at $\abz n\le2^{18}$, tail $\frac{2\abz J}\pi N^{-1/2}$
 & $18.70397$; increment $2^{16}\!\to\!2^{18}$ is $0.00124$ vs predicted $\frac2\pi(\frac1{256}-\frac1{512})=0.001243$
 & ok\\
(d) partial sums $481.6,644.6$ at $N=2^{16},2^{18}$; increment $162.97$
 & $481.648$, $644.619$; increment $162.971$ vs $\frac2\pi(512-256)=162.975$ & ok\\
(f) $\sum_{n\ne0}\abz n(n^2-1)\abz{\hat G_n}^2=4.076764$, $\nrm{T(\gtot)\Om}^2=0.169865\,c$
 & $4.07676431$ (converged at $2^{10}$); $/24=0.1698652$ & ok\\
(g) $\nrm{L_0^{1/2}T(\gtot)\Om}^2=1.622737\,c$; $L_0^2$-moment $+0.1405$ per factor $16$
 & $38.945685/24=1.6227369$; increments $0.14044$, $0.14046$, then $0.07411$ for a factor $4$ & ok\\
(e) $H^s$ sums converge for $s=2,2.4$, log-diverge at $2.5$, power-diverge at $2.6$
 & $41.714334$; $362.2476$ (slow); $+0.1404$ per factor 16; $+0.75,+1.31$ growing & ok\\
\bottomrule
\end{tabularx}
\end{center}
The script is a faithful implementation of the analytic claims: the field is built in the $x$ chart with the
smooth bump cut-off, pushed to the circle by \eqref{eq:phi} below, and Fourier-transformed by FFT with the
correct $(-1)^n$ phase for the grid $\theta_j=-\pi+2\pi j/N$.  I checked the two places where such scripts
usually go wrong: the point at infinity ($\abz{\cos(\theta/2)}<10^{-14}$ is mapped to $G=0$, correct because
$G(\theta)=2\cos^2(\theta/2)g(\tan(\theta/2))$ and $g$ is compactly supported in $x$ for representative (i)),
and the normalisation
\begin{equation}\label{eq:phi}
 \hat G_n=\frac1{2\pi}\int_{-\pi}^\pi G e^{-in\theta}\dd\theta\approx(-1)^n\,\frac{\mathrm{fft}(G)_n}N ,
\end{equation}
which is right.  The only caveat is that the reported convergence of the $s=2.4$ Sobolev sum is slow
($N^{-0.2}$) and is presented as ``converges''; it does, but the numbers alone do not distinguish $s=2.4$ from
a very slow divergence --- the analytic argument (Theorem~2.2(b),(e)) is what settles it.

\section{Overall verdict on Theorem~5.1}\label{sec:verdict}

\begin{verdictok}
\textbf{Theorem~5.1 stands, for every diffeomorphism covariant conformal net on $S^1$ in the sense of CDIT \S4
(axioms (1)--(10), separable $\Hil$), with arbitrary central charge $c>0$, no strong additivity, no finiteness
of $L_0$-multiplicities, no rationality} --- subject to the single write-up gap of \S\ref{sec:phase} (the phase
in $U(\Exp f)=e^{iT(f)}$ must be scalar on a reducible $\Diff$-representation; standard, repairable by citing
the global lift to the Virasoro group, and used in the same generality by CW05 \S5).  I could not find a
sub-class restriction hiding anywhere: the three external inputs (X1)--(X3) are, respectively, CW05 Thm.~4.4
(proved for a general positive-energy representation, constant independent of the lowest weight --- I read the
proof), CW05 Prop.~4.5/Lemma~4.6 (likewise), and the automatic properties (6)--(9) of a M\"obius covariant net.
The genuinely new content of the note --- and it \emph{is} new relative to both sources --- is:
(1) the observation that CDIT's $s>3$ is a restriction on the group-level extension and is bypassed for a
one-parameter flow by Stone's theorem, which I confirm is not circular (nothing in CW05 Thm.~4.4 or
Prop.~4.5 requires the diffeomorphism, only the field); (2) the split $\gtot=(\gtot-m)+m$, which is what makes
the Trotter/localisation argument applicable to a field that vanishes on one side and is M\"obius on the other
(\S\ref{sec:trotter}); and (3) the extension of the covariance statement to intervals straddling the touching
point via CDIT Prop.~4.1's approximation argument (\S\ref{sec:straddle}).
\end{verdictok}

\paragraph{Consequences for the Phase-2 chain.}  With (H1)--(H3) unconditional, Theorem~B of
\texttt{universality\_theorem\_B.tex} (which additionally needs only (HD), (BW), (RS), all automatic for
M\"obius covariant nets) becomes unconditional as well, and with it $f_2=c/(12\pi^2)$ for every diffeomorphism
covariant conformal net.  I checked the interface: Theorem~B consumes $a=T(\gtot)\Om\in\Hil$ (supplied by
Theorem~2.2(f)), $\varphi(K)=-2\operatorname{Im}\ipp a{K\Om}$ on $\cA(I)_{\mathrm{sa}}$ (supplied by
\S\ref{sec:H3} in exactly that form), $\omega_s(K^2)\to\omega(K^2)$ (supplied), and the unitaries $U_s$ with
$\Ad U_s|_{\cA(A)}=\mathrm{id}$, $\Ad U_s|_{\cA(D)}=\Ad U(h_s)$ (supplied for $s\ge0$, which is the sign used).
Nothing in Theorem~B needs $T(\gtot)\Om\in D(L_0)$, so the logarithmic failure of that (Theorem~2.2(g)) and the
now-\emph{closed} form-bound route (\S\ref{sec:form}) do not touch the main line.  What remains genuinely open
after this report, and should be stated as such in the note, is only:
\begin{enumerate}[leftmargin=1.6em,label=\textup{(O\arabic*)}]
\item $\ipp{U_s^*\Om}{L_0U_s^*\Om}<\infty$ for fixed $s>0$ (Note~6 Conj.~3.2).  Route~B (the Schwarzian energy
      integral, note \S6 (B2)) is now the \emph{only} candidate, since the commutator route is dead
      (\S\ref{sec:form}); the missing step there is the weak-$*$ convergence of the Schwarzian measures of the
      mollified approximants, which the note flags.
\item the group-level extension of $U$ to $D^\sigma(S^1)$, $2<\sigma\le3$, for general $c$ (open in CDIT;
      known for integer $c$ in Fock space via [DIT19]).  Needed only for compositions of maps with different
      touching points, not for Phase~1 or Phase~2.
\end{enumerate}

\paragraph{Recommended edits to \texttt{implementation\_C11.tex}.}
(1) Replace Remark~5.2 and the last table row by the sharp statement of \S\ref{sec:form}
($\beta=1$ true and sharp; $\nrm{\gtot'}_1=\infty$; route closed).
(2) Add one sentence and a citation for the scalar phase (\S\ref{sec:phase}).
(3) Insert the half-page of \S\ref{sec:straddle}$\langle1\rangle2$ (the union of the $\varphi^{(s)}_kJ$ is an
interval decreasing to $\overline{\ktil_sJ}$).
(4) Fix (M1), (M2), (M3), (M5), (M6); state (M4)'s $s\ge0$.
(5) Optional but worth it: record that CW05 \S4 is already stated for reducible representations, so Lemma~4.2 is
needed only to use CDIT's phrasing.

\section{Summary of findings}\label{sec:summary}

\begin{center}\small
\begin{tabularx}{\textwidth}{l >{\raggedright\arraybackslash}X}\toprule
Severity & Finding\\\midrule
BREAK & \S\ref{sec:form}: Remark~5.2 / last table row.  ``Form bound with $\beta<1$, window
$\frac12\le\beta<1$ open'' is false; the optimally weighted test vectors force $\beta\ge1$, $\beta=1$ holds, and
$\nrm{\gtot'}_1=\infty$.  Does not affect (H1)--(H3).\\
GAP & \S\ref{sec:phase}: $U(\Exp f)=e^{iT(f)}$ modulo a \emph{scalar} phase is cited from CDIT \S2.1, which
states it on the irreducible $\Hil(c,h)$; on a reducible $\Diff$-representation the a priori discrepancy is
$\bigoplus_je^{i\alpha_j}$.  Used essentially in the localisation step, hence in all of (H1).\\
MINOR & \S\ref{sec:straddle}: ``for $m$ large the union is an interval'' is asserted, not proved.\\
MINOR & \S\ref{sec:H3}: $\varphi(X)=i\omega([T(\gtot),X])$ is formal ($X\Om\in D(\overline{T(\gtot)})$ unknown);
the sesquilinear expression is the rigorous one and is what Theorem~B uses.\\
MINOR & (M1) factor $4$ in Prop.~5.1(ii); (M2) sign of $[L_0,T(g)]$; (M3),(M4) the $s\ge0$ restrictions;
(M5) one-sided vs two-sided sums in the numerics box; (M6) separability of $\Hil$.\\
OK & Everything else: Thm.~2.2 (a)--(g), Prop.~2.1, Lemma~4.2, Thm.~4.4 (i)--(iii), Thm.~4.5 (H2),(H3),
Prop.~5.1 (i),(iii),(iv), \S6 (B1)--(B3), and all numerical claims.\\
\bottomrule
\end{tabularx}
\end{center}

\begin{thebibliography}{9}
\bibitem{CW05} S.~Carpi, M.~Weiner, \emph{On the uniqueness of diffeomorphism symmetry in conformal field
theory}, CMP \textbf{258} (2005) 203, arXiv:math/0407190; \texttt{refs/P1\_CarpiWeiner\_math-0407190.pdf}.
Read for this report: pp.~11--18 (Lemma 4.1, Props.~4.2, 4.3, Thm.~4.4, Prop.~4.5, Lemma~4.6, Lemma~5.3,
Props.~5.1, 5.2, 5.4, Thm.~5.5), including all proofs.
\bibitem{CDIT} S.~Carpi, S.~Del Vecchio, S.~Iovieno, Y.~Tanimoto, \emph{Positive energy representations of
Sobolev diffeomorphism groups of the circle}, AHP \textbf{22} (2021) 2091, arXiv:1808.02384;
\texttt{refs/CDIT\_1808.02384.pdf}.  Read: pp.~5--6, 22--24.
\bibitem{FH05} C.~J.~Fewster, S.~Hollands, arXiv:math-ph/0412028;
\texttt{refs/FewsterHollands05\_math-ph-0412028.pdf}.  Used only through CDIT Prop.~2.1 and for the
Virasoro-group lift suggested in \S\ref{sec:phase} (their Prop.~5.1); \textbf{not re-read in detail here}.
\bibitem{RS} M.~Reed, B.~Simon, \emph{Methods of Modern Mathematical Physics I--II}.  \textbf{NOT OBTAINED};
Thms.~VIII.21, VIII.25(a), VIII.31, X.37 used as standard, as in the note.
\bibitem{QB} Track QB, \texttt{rigor/implementation\_C11.tex} (2026-09-08), the document under review;
\texttt{rigor/qb\_regularity.py}, re-run for \S\ref{sec:numerics}.
\bibitem{QA} Track QA, \texttt{rigor/universality\_theorem\_B.tex} (2026-09-08), for the interface check in
\S\ref{sec:verdict}.
\end{thebibliography}
\end{document}
